\section*{5 Axis-evasive partitions: 5.2 Explicit constructions}
In this section we give explicit constructions of axis-evasive partitions (see Definition35). Our constructions are based on quadratic field extensions. We identify a set $P$ of size $q$--a prime power--with the finite field $\mathbb F_q$ in an arbitrary manner, namely, by using an arbitrary bijection which, for ease of readability, we do not make explicit in the notation. We start by giving some basic background on finite fields.
Let $h(x)\in\mathbb F_q[x]$ be a degree $2$ irreducible monic polynomial. It is a well-known fact that $\mathbb F_q[x]/\langle h(x)\rangle$ is a field of size $q^2$ which we denote, somewhat less informatively, by $\mathbb F_{q^2}$. Note that there exists $\alpha\in\mathbb F_{q^2}$ such that $h(\alpha)=0$ (indeed, take $\alpha=x+\langle h(x)\rangle$). Since $h$ is irreducible over $\mathbb F_q$ and has degree $2$, we can write every element of $\mathbb F_{q^2}$ in the form $a+\alpha b$, where $a,b\in\mathbb F_q$, in a unique manner. That is, we can identify in the set-theoretic level, $\mathbb F_{q^2}$ with $\mathbb F_q+\alpha\mathbb F_q$. Using this identification, we identify $P^2$ with $\mathbb F_{q^2}$ in the natural way, namely, a point $(a,b)\in P^2$ is identified with $a+\alpha b$ in $\mathbb F_{q^2}$. Note that, with this identification, the horizontal lines in $P^2$ are of the form $b\alpha+\mathbb F_q$ where $b\in\mathbb F_q$ can be thought of as the fixed height of the line. Similarly, the vertical lines are given by $b+\alpha\mathbb F_q$. Given $\delta\in\mathbb F_{q^2}\setminus\{0\}$, we say that $\ell_\delta=\delta\mathbb F_q\subseteq\mathbb F_{q^2}$ is the line through the origin with slope $\delta$.
Our construction of exis-evasive partitions is based on an equivalence relation that we are about to define. The partition is then obtained by considering the respective equivalence classes. We begin the construction by ignoring the ``origin'' $0\in\mathbb F_{q^2}$ and work only with $\mathbb F_{q^2}\setminus\{0\}$. Note that this is the set of invertible elements of $\mathbb F_{q^2}$ which has a group structure under the field multiplication. When referring to this multiplicative group we write $(\mathbb F_{q^2})^\times$.
Let $\beta\in(\mathbb F_{q^2})^\times$. Denote by $o(\beta)$ the order of $\beta$ in the multiplicative group $(\mathbb F_{q^2})^\times$. It will be convenient to denote $c=o(\beta)-1$. We define an equivalence relation on $(\mathbb F_{q^2})^\times$, parameterized by $\beta$, as follows: For $\gamma,\delta\in(\mathbb F_{q^2})^\times$
\begin{equation}\gamma\sim\delta\quad\Longleftrightarrow\quad\gamma\delta^{-1}\in\langle\beta\rangle,\tag{5.1}\end{equation}
where $\langle\beta\rangle$ is the subgroup of $(\mathbb F_{q^2})^\times$ that is generated by $\beta$. Observe that this is an equivalence relation. Indeed, the classes are the different cosets, that is, the elements of the quotient group $(\mathbb F_{q^2})^\times/\langle\beta\rangle$. For completeness, we quickly prove that this is an equivalence relation: as $1\in\langle\beta\rangle$, we have that $\gamma\sim\gamma$. Secondly, if $\gamma\delta^{-1}\in\langle\beta\rangle$ then $\delta\gamma^{-1}\in\langle\beta^{-1}\rangle=\langle\beta\rangle$ which establishes symmetry. As for transitivity, if $\gamma\sim\delta$ and $\delta\sim\varepsilon$ then
\[\gamma\varepsilon^{-1}=\gamma(\delta^{-1}\delta)\varepsilon^{-1}=(\gamma\delta^{-1})(\delta\varepsilon^{-1})\in\langle\beta\rangle.\]
One can easily see that the equivalence class of an element $\gamma\in(\mathbb F_{q^2})^\times$ is $[\gamma]=\gamma\langle\beta\rangle=\{\gamma,\beta\gamma,\ldots,\beta^c\gamma\}$. Note further that $|[\gamma]|=c+1$. Indeed, if there are $0\leq j<i\leq c$ such that $\beta^i\gamma=\beta^j\gamma$ then $0=(\beta^i-\beta^j)\gamma=(\beta^{i-j}-1)\beta^j\gamma$, which is a contradiction as none of the factors in the product is zero.
In the following claim we show that, under some conditions on $\alpha,\beta$, the second property of axis-evasiveness is met by the construction above. We mention already here that the third condition in Definition35 is not met by the construction as is (regardless of the choice of $\alpha,\beta$), and we will alter it afterwards to meet that property as well.
\setcounter{nttheorem}{56}
\begin{ntclaim}[Author Claim57]
Assume that $\langle\beta\rangle\cap\ell_\alpha=\langle\beta\rangle\cap\ell_{\alpha^{-1}}=\varnothing$ and that $\langle\beta\rangle\cap\mathbb F_q=\{1\}$. Then, for every $\ell,\ell'\in\mathcal X$ (not necessarily distinct) it holds that $\ell'\cap(\ell)\leq c$.
\end{ntclaim}
\begin{proof}
Recall that $(\gamma)=\{\beta\gamma,\ldots,\beta^c\gamma\}$. Thus,
\[\bigcup_{\gamma\in\ell}(\gamma)=\bigcup_{\gamma\in\ell}\bigcup_{i=1}^c\{\beta^i\gamma\}=\bigcup_{i=1}^c\beta^i\ell.\]
Therefore,
\begin{equation}\ell'\cap(\ell)=\ell'\cap\bigcup_{\gamma\in\ell}(\gamma)=\bigcup_{i=1}^c(\ell'\cap\beta^i\ell).\tag{5.2}\end{equation}
Fix $i\in[c]$ and consider two cases. First, if $\ell$ is vertical, namely, $\ell=b+\alpha\mathbb F_q$ for some $b\in\mathbb F_q$, then $\beta^i\ell=\beta^ib+\alpha\beta^i\mathbb F_q$. Since, by assumption, $\langle\beta\rangle\cap\mathbb F_q=\{1\}$ we have that $\alpha\beta^i\mathbb F_q\ne\alpha\mathbb F_q$ and so the line $\beta^i\ell$ is not vertical. As, by assumption, $\langle\beta\rangle\cap\ell_{\alpha^{-1}}=\varnothing$, we have that $\alpha\beta^i\notin\mathbb F_q$ and so the line $\beta^i\ell$ is not horizontal either.
Second, consider the case that $\ell$ is horizontal $\ell=b\alpha+\mathbb F_q$ for some $b\in\mathbb F_q$. Then, $\beta^i\ell=b\alpha\beta^i+\beta^i\mathbb F_q$. Per our assumption that $\langle\beta\rangle\cap\ell_\alpha=\varnothing$, we have that $\beta^i\mathbb F_q\ne\alpha\mathbb F_q$ and so the line $\beta^i\ell$ is not vertical. As we assume $\langle\beta\rangle\cap\mathbb F_q=\{1\}$, we have that $\beta^i\mathbb F_q\ne\mathbb F_q$, and so the line $\beta^i\ell$ cannot be horizontal either. To summarize, we have that $\beta^i\ell\notin\mathcal X$. However, $\ell'\in\mathcal X$ and so $\beta^i\ell$ and $\ell'$ are two distinct lines. As such, the two lines intersect in at most one point. Equation(5.2) then yield $|\ell'\cap(\ell)|\leq c$.
\end{proof}
\paragraph{Informal discussion regarding the third property}
As mentioned above, the partition of $(\mathbb F_{q^2})^\times$ as defined above does not have the third property required for axis-evasiveness. Namely, there are $\gamma\in(\mathbb F_{q^2})^\times$ such that $[\gamma]$ intersects some axis-parallel line at more than one point. To get some idea on which equivalence classes $[\gamma]$ are problematic, let us first ask when do $\gamma,\beta\gamma$ are on some common axis-parallel line. We first observe that two points $\delta,\varepsilon\in(\mathbb F_{q^2})^\times$ are on a common axis-parallel line if and only if $\delta-\varepsilon\in\{1,\alpha\}\mathbb F_q$. Thus, $\gamma$ and $\beta\gamma$ are on the same axis-parallel line if and only if $\gamma-\beta\gamma=(1-\beta)\gamma\in\{1,\alpha\}\mathbb F_q$. This is equivalent to saying that $\gamma$ is on one of the two lines through the origin with slopes $1/(1-\beta),\alpha/(1-\beta)$.
More generally, $[\gamma]$ intersects with some axis-parallel line in more than one point if and only if $\beta^i\gamma-\beta^j\gamma\in\{1,\alpha\}\mathbb F_q$ for some $0\leq j<i\leq c$. Equivalently, $\gamma$ is on a line $\ell_\delta$ with
\begin{equation}\delta\in\left\{\frac1{\beta^i-\beta^j},\frac\alpha{\beta^i-\beta^j}\;\middle|\;0\leq j<i\leq c\right\}.\tag{5.3}\end{equation}
The key observation is that although there are a fair amount of ``bad'' points $\gamma$, they are all contained in a small number of lines. By ``small'' here we mean that the number is polynomial in $c$ and is independent of $q$. Thus, the hope is that by redefining the partition on these few problematic lines we will not harm the previous analysis by much. Indeed, no matter how we alter the partition restricted to these lines, if we make sure none of them is axis-parallel (by requiring more properties from $\alpha,\beta$) then each of these lines intersect an axis-parallel line at one point. As a result, the bound obtained in Claim57 will deteriorate proportionally to the number of lines above.
The only small technical issue is that even if $\gamma\in\ell_\delta$ for some slope $\delta$ as above, it is not the case that $[\gamma]\subseteq\bigcup_\varepsilon\ell_\varepsilon$ where $\varepsilon$ is taken from the set of slopes given by Equation(5.3). As we wish to alter the partition defined above, it would be cleaner to have all of the points in $[\gamma]$ of a problematic point $\gamma$ contained in the set of points on which we redefine the partition. Thus, we ``close'' the set of slopes given by Equation(5.3) to multiplication by $\beta$.
Ending the informal discussion and returning to the formal analysis, we consider the set of slopes.
\begin{equation}\Delta=\left\{\frac{\beta^k}{\beta^i-\beta^j},\frac{\alpha\beta^k}{\beta^i-\beta^j}\;\middle|\;0\leq j<i\leq c\text{ and }0\leq k\leq c\right\}.\tag{5.4}\end{equation}
Further define the set of all points in $(\mathbb F_{q^2})^\times$ covered by the lines with slopes from $\Delta$ by $U=\bigcup_{\delta\in\Delta}\ell_\delta$. This definition of $\Delta$ indeed fixes the technical caveat discussed above, as the following claim states.
\begin{ntclaim}[Author Claim58]
For every $\gamma\in(\mathbb F_{q^2})^\times$ either $[\gamma]\subseteq U$ or $[\gamma]\cap U=\varnothing$.
\end{ntclaim}
\begin{proof}
If an element $\varepsilon\in U$ then $\varepsilon\in\ell_\delta$ for some $\delta\in\Delta$. Note that $\beta\varepsilon\in\ell_{\beta\delta}$ and that $\beta\delta\in\Delta$. Hence, $\beta\varepsilon\in U$. Therefore, $\varepsilon\in U\Longrightarrow\varepsilon\langle\beta\rangle\subseteq U$. Assume now that $[\gamma]\cap U\ne\varnothing$, and take $\gamma\beta^i\in U$. By the above, $\gamma\beta^i\langle\beta\rangle\subseteq U$. The proof then follows as $\gamma\beta^i\langle\beta\rangle=\gamma\langle\beta\rangle=[\gamma]$.
\end{proof}
Define a new partition of $\mathbb F_{q^2}$ (including $0$) which agrees with the one that is given by Equation(5.1) on $\mathbb F_{q^2}^\times\setminus U$. By Claim58, this is well-defined. The new partition, restricted to $U$, is done as follows. Let $\delta_0\in\Delta$ be an arbitrary element. Note that
\[U=\ell_{\delta_0}\cup\bigcup_{\delta\in\Delta\setminus\{\delta_0\}}(\ell_\delta\setminus\{0\})\]
is a disjoint union. To partition $U$, we partition $\ell_{\delta_0}$ as well as each of $\ell_\delta\setminus\{0\}$ where $\delta\in\Delta\setminus\{\delta_0\}$ in an arbitrary way provided it has the least number of parts under the conditions that each part has size at most $c+1$. For ease of readability, we denote by $[\gamma]$ the class with respect to the new partition.
\begin{ntclaim}[Author Claim59]
Assume, on top of the assumptions of Claim57 that for every $\delta\in\Delta$, $\ell_\delta\notin\mathcal X$. Then, the new partition defined above is $(c,4c^2)$-axis evasive.
\end{ntclaim}
\begin{proof}
First, observe that by construction, every class intersects any axis-parallel line in at most one point. Indeed, classes that are outside of $U$ have this property by the definition of $U$ as can be easily verified (and discussed above). Moreover, by the way we redefined the partition restricted to $U$, every class that is a subset of $U$ is also a subset of a line $\ell_\delta$ for some $\delta\in\Delta$. As $\ell_\delta\notin\mathcal X$ by hypothesis, we have that the line and, as a result, the class it contains, intersects any axis-parallel line in at most one point. This establishes the third property of axis-evasiveness. The second property follows as, by construction, every part has size at most $c+1$.
Moving on to the second property, consider $\ell,\ell'\in\mathcal X$, not necessarily distinct. As outside of $U$ the partition is defined as before, Claim58 yields
\begin{equation}\left|\ell'\cap\bigcup_{\gamma\in\ell\setminus U}(\gamma)\right|\leq c.\tag{5.5}\end{equation}
Take $\gamma\in U\cap\ell$. Since, by construction $(\gamma)\subseteq\ell_\delta$ for some $\delta\in\Delta$, and since by hypothesis $\ell_\delta\notin\mathcal X$ we have that $|\ell'\cap\ell_\delta|=1$ and $(\gamma)\cap\ell'\subseteq\ell_\delta\cap\ell'$. Therefore, $|(\gamma)\cap\ell'|\leq1$. Together with Equation(5.5) we get that $|\ell'\cap(\ell)|\leq c+|U\cap\ell|$. Now, since $\ell\in\mathcal X$ and every line $\ell_\delta$ with slope $\delta\in\Delta$ is not in $\mathcal X$ we have that $|\ell\cap\ell_\delta|=1$. Thus, $|U\cap\ell|\leq|\Delta|$ which implies $|\ell'\cap(\ell)|\leq c+|\Delta|$.
To conclude the proof, we turn to bound $|\Delta|$. It is straightforward to give a bound of $O(c^3)$ though one can optimize the bound a bit. Indeed, with the notation of Equation(5.4), by multiplying by $\beta^{-\min(j,k)}$, one can rewrite
\begin{align}\Delta&=\left\{\frac1{\beta^i-\beta^j},\frac\alpha{\beta^i-\beta^j}\;\middle|\;0<j<i\leq c\right\}\notag\\&\quad\cup\left\{\frac{\beta^j}{\beta^i-1},\frac{\alpha\beta^j}{\beta^i-1}\;\middle|\;0<i\leq c,\ 0\leq j\leq c\right\}.\tag{5.6}\end{align}
Thus, $|\Delta|\leq3c^2$, and the proof follows.
\end{proof}
We summarize the discussion so far.
\begin{ntproposition}[Author Proposition60]
Let $\mathbb F_q$ be finite field. Let $h(x)\in\mathbb F_q[x]$ be a degree $2$ irreducible monic polynomial, and consider the field $\mathbb F_q[x]/\langle h(x)\rangle$ which we denote by $\mathbb F_{q^2}$. Let $\alpha,\beta\in\mathbb F_{q^2}$ be two elements satisfying:
\begin{enumerate}
\item $h(\alpha)=0$,
\item $\langle\beta\rangle\cap\mathbb F_q=\{1\}$,
\item $c+1=o(\beta)\leq\sqrt q/10$,
\item $\langle\beta\rangle\cap\ell_\alpha=\langle\beta\rangle\cap\ell_{\alpha^{-1}}=\varnothing$,
\item $(\langle\beta\rangle-\langle\beta\rangle)\cap\mathbb F_q=\{0\}$,
\item $(\langle\beta\rangle-\langle\beta\rangle)\cap\ell_\alpha=(\langle\beta\rangle-\langle\beta\rangle)\cap\ell_{\alpha^{-1}}=\{0\}$.
\end{enumerate}
Then, there exists a partition of $(\mathbb F_q)^2$ that is $(c,4c^2)$-axis-evasive, where $c=o(\beta)-1$. The number of parts in the partition is bounded above by $2q^2/(c+1)$.
\end{ntproposition}
To prove Proposition60 we need the following easy claim.
\begin{ntclaim}[Author Claim61]
Let $\delta\in(\mathbb F_{q^2})^\times$ be such that $\delta\notin\mathbb F_q\cup\ell_\alpha$ then, $\ell_\delta\notin\mathcal X$.
\end{ntclaim}
\begin{proof}
Write $\delta=a+\alpha b$ with $a,b\in\mathbb F_q$. Then, $\ell_\delta=(a+\alpha b)\mathbb F_q$. Observe that if $\ell_\delta$ is vertical then $a=0$ and so $\delta\in\ell_\alpha$. Similarly, if $\ell_\delta$ is horizontal then $b=0$ implying $\delta\in\mathbb F_q$.
\end{proof}
\begin{proof}[Author proof of Proposition60]
To bound the number of parts, recall that in the original partition, each part has size $c+1$. Moreover, in the altered partition we partition each line $\ell_\delta$ with slope $\delta\in\Delta$ (excluding the origin from all but for one of the lines $\ell_{\delta_0}$) to parts of size $c+1$ each, except for possibly one part. As $|\Delta|\leq3c^2$, the number of parts it bounded by
\[\frac{q^2-1}{c+1}+|\Delta|\left(1+\frac q{c+1}\right)\leq\frac{q^2-1}{c+1}+6cq\leq\frac{2q^2}c,\]
where the last inequality follows by our assumption that $o(\beta)\leq\sqrt q/10$.
To conclude the proof of the proposition, it suffices to show that for every $\delta\in\Delta$ it holds that $\ell_\delta\notin\mathcal X$. By Claim61, it suffices to prove that $\delta\notin\mathbb F_q\cup\ell_\alpha=\{1,\alpha\}\mathbb F_q$. There are two types of slopes $\delta\in\Delta$, according to whether they appear in the first or second set in Equation(5.6). The first kind is of the form
\[\delta=\frac{\alpha^k}{\beta^i-\beta^j},\]
with $0<j<i\leq c$ and $k\in\{0,1\}$. If $\delta\in\{1,\alpha\}\mathbb F_q$ then $\delta^{-1}\in\{1,\alpha^{-1}\}\mathbb F_q$ and so $\beta^i-\beta^j\in\{\alpha^k,\alpha^{k-1}\}\mathbb F_q$ in contradiction to our hypothesis. Consider now the other kind of slope
\[\delta=\frac{\alpha^k\beta^j}{\beta^i-1}\]
where $0<i\leq c$, $0\leq j\leq c$ and $k\in\{0,1\}$. If $\delta\in\{1,\alpha\}\mathbb F_q$ then $\delta^{-1}\in\{1,\alpha^{-1}\}\mathbb F_q$ and so $(\beta^i-1)\beta^{-j}\in\{\alpha^k,\alpha^{k-1}\}\mathbb F_q$. Note that $(\beta^i-1)\beta^{-j}=\beta^{i-j}-\beta^{-j}\in\langle\beta\rangle-\langle\beta\rangle$ and so we again get a contradiction.
\end{proof}
We are now ready to prove Theorem47. For the sake of readability, we repeat its statement here.
\begin{nttheorem}[Author Theorem62]
Let $P$ be a set of size $q$, where $q$ is an odd prime power. Let $c$ be an even integer such that $c+1\mid q+1$, and $c\leq\sqrt q/10$. Then, there exists a $(c,4c^2)$-axis evasive partition of $P^2$ with at most $2q^2/(c+1)$ parts.
\end{nttheorem}
\begin{proof}
As above, we identify $P^2$ with $\mathbb F_{q^2}$. It is a well-known fact that the multiplicative group $(\mathbb F_{q^2})^\times$ is cyclic. A basic result in group theory states that a cyclic group has a (unique) subgroup of every given size which divides the group size. Now, $|(\mathbb F_{q^2})^\times|=q^2-1=(q-1)(q+1)$. Thus, as $c+1\mid q+1$, there exists a subgroup $H$ of $(\mathbb F_{q^2})^\times$ of size $c+1$. The subgroup $H$ is cyclic, being a subgroup of a cyclic group. Let $\beta$ be a generator for $H$. We first prove that $\beta$ satisfies those hypothesis of Proposition60 that do not involve $\alpha$, namely, conditions(2) and(5).
\begin{ntclaim}[Author Claim63]
$(\langle\beta\rangle-\langle\beta\rangle)\cap\mathbb F_q=\{0\}$ and $\langle\beta\rangle\cap\mathbb F_q=\{1\}$.
\end{ntclaim}
\begin{proof}
Assume towards a contradiction that $\beta^i-\beta^j\in\mathbb F_q$ for some $0\leq j<i\leq c$. Since $x^q=x$ for every $x\in\mathbb F_q$, we get
\[\beta^i-\beta^j=(\beta^i-\beta^j)^q=\beta^{iq}-\beta^{jq},\]
where the last equality follows since $q$ is divisible by the characteristic of the field. Recall that $o(\beta)=c+1\mid q+1$ and so $\beta^{i(q+1)}=1$, implying $\beta^{iq}=\beta^{-i}$. Thus,
\[\beta^i-\beta^j=\frac1{\beta^i}-\frac1{\beta^j}=\frac{\beta^j-\beta^i}{\beta^{i+j}}.\]
As $\beta^i\ne\beta^j$ the above equation implies $\beta^{i+j}=-1$, and so $-1\in H$. Since $q$ is odd, the characteristic of the field $\mathbb F_{q^2}$ is odd and so $o(-1)=2$. Lagrange's Theorem then implies that $2\mid|H|=c+1$, which stands in contradiction to $c$ being even.
To prove that $\langle\beta\rangle\cap\mathbb F_q=\{1\}$, take $\beta^i$ with $0<i\leq c$. If $\beta^i\in\mathbb F_q$ then $\beta^{iq}=\beta^i$. On the other hand, we proved above that $\beta^{iq}=\beta^{-i}$, and so $\beta^i=\beta^{-i}$ implying $\beta^{2i}=1$. Therefore, $o(\beta)=c+1\mid2i$, but this is impossible as $0<i\leq c$ and, recall, $c$ is even.
\end{proof}
We proceed with the proof of Theorem47 by finding $\alpha\in\mathbb F_{q^2}$ that, together with the already chosen $\beta$, satisfies the remaining conditions in the hypothesis of Proposition60. Since $\mathbb F_{q^2}$ is a quadratic field extension of $\mathbb F_q$, every element $\gamma\in\mathbb F_{q^2}\setminus\mathbb F_q$ has degree $2$. That is, the minimal polynomial $h_\gamma$ of every such $\gamma$ over $\mathbb F_q$ is of degree $2$ (and can be made monic by dividing by the leading coefficient, if necessary). Indeed, $\deg(h_\gamma)$ cannot equal $1$ as this would imply $\gamma\in\mathbb F_q$. On the other hand,
\[2=[\mathbb F_{q^2}:\mathbb F_q]=[\mathbb F_{q^2}:\mathbb F_q(\gamma)][\mathbb F_q(\gamma):\mathbb F_q]=[\mathbb F_{q^2}:\mathbb F_q(\gamma)]\deg(h_\gamma),\]
which shows that if $\deg(h_\gamma)\ne1$ then $\deg(h_\gamma)=2$.
Thus, condition(1) in the hypothesis of Proposition60 holds for every element in $\mathbb F_{q^2}\setminus\mathbb F_q$. Hence, to prove that all the remaining conditions in the hypothesis of Proposition60 hold, it suffices to prove that there exists $\alpha\in\mathbb F_{q^2}\setminus\mathbb F_q$ which satisfies conditions(4) and(6). To this end, pick a set of ``slopes'' $\Delta'=\{\delta_1,\ldots,\delta_{q+1}\}\subseteq(\mathbb F_{q^2})^\times$ such that $(\mathbb F_{q^2})^\times$ is the disjoint union
\[(\mathbb F_{q^2})^\times=\bigcup_{\delta\in\Delta'}(\ell_\delta\setminus\{0\}).\]
For example, $\Delta'=\{a+\alpha\mid a\in\mathbb F_q\}\cup\{1\}$ will do. For $\delta\in(\mathbb F_{q^2})^\times$ let
\[I_\delta=|\langle\beta\rangle\cap\ell_\delta|+|(\langle\beta\rangle-\langle\beta\rangle)\cap(\ell_\delta\setminus\{0\})|.\]
Since the $\ell_\delta\setminus\{0\}$ with $\delta\in\Delta'$ are disjoint, $0\notin\langle\beta\rangle$, and since $|\langle\beta\rangle|=c+1$ and $|\langle\beta\rangle-\langle\beta\rangle|\leq(c+1)^2$, we have that
\[\mathbb E_\delta[I_\delta]\leq\frac{(c+1)^2+(c+1)}{q+1}\leq\frac{2(c+1)^2}{q+1},\]
where $\delta$ is sampled uniformly from $\Delta'$. By Markov's inequality, for at least $3/4$ of the elements $\delta\in\Delta'$ it holds that
\[|I_\delta|\leq\frac{8(c+1)^2}{q+1}.\]
Note that $(\mathbb F_{q^2})^\times$ is also a disjoint union of $\{\ell_{\delta^{-1}}\setminus\{0\}\mid\delta\in\Delta\}$. Thus, using the same argument as above, we get that for at least $1/2$ the elements $\delta\in\Delta'$, both $|I_\delta|$ and $|I_{\delta^{-1}}|$ are bounded by $8(c+1)^2/(q+1)$. But, as $c\leq\sqrt q/10$, this bound is strictly smaller than $1$, implying that $I_i=I_{q+1-i}=0$. That is, at least half the elements $\delta\in\Delta'$ satisfy conditions(4) and(6). Take $\alpha$ to be any of these elements. To conclude, we found $\alpha$ and $\beta$ for which all the conditions in the hypothesis of Proposition60 are met, and the proof follows.
\end{proof}
